\documentclass[11pt]{article}

\usepackage[margin=0.9in]{geometry}
\usepackage{amsmath,amssymb}
\usepackage{booktabs}
\usepackage{enumitem}
\usepackage{listings}
\usepackage{xcolor}
\usepackage{hyperref}

\definecolor{backcolour}{rgb}{0.95,0.95,0.92}
\definecolor{codeblue}{rgb}{0.05,0.15,0.55}
\definecolor{codegreen}{rgb}{0.0,0.45,0.0}

\lstset{
    backgroundcolor=\color{backcolour},
    basicstyle=\ttfamily\small,
    keywordstyle=\color{codeblue}\bfseries,
    commentstyle=\color{codegreen},
    stringstyle=\color{purple},
    breaklines=true,
    frame=single,
    showstringspaces=false,
    tabsize=4
}

\title{CPSC 250 \\ Plotting Error Bands}
\author{Edward Brash}
\date{}

\begin{document}

\maketitle

\section{Motivation}

When we fit a curve to data, the fit parameters have uncertainties.

A natural question is:

\begin{quote}
How do uncertainties in the fit parameters translate into uncertainties in the fit itself?
\end{quote}

The full statistical treatment is complicated.

However, we can create a useful visualization called an \textbf{error band}.

\section{The Basic Idea}

Suppose our fit parameters are stored in:

\begin{itemize}
    \item \texttt{popt} --- optimal fit parameters,
    \item \texttt{pcov} --- covariance matrix.
\end{itemize}

The covariance matrix describes the uncertainties and correlations of the parameters.

We can use this information to generate many random possible fits.

\section{Monte Carlo Sampling}

The key idea is:

\begin{enumerate}
    \item randomly generate many possible parameter sets,
    \item compute a fit curve for each parameter set,
    \item examine the spread of the curves.
\end{enumerate}

\section{Generating Random Parameter Sets}

We use:

\begin{lstlisting}[language=Python]
ps = np.random.multivariate_normal(
    popt,
    pcov,
    10000
)
\end{lstlisting}

This generates:

\begin{itemize}
    \item 10,000 random parameter sets,
    \item sampled from a multivariate Gaussian distribution,
    \item centered on the best-fit parameters.
\end{itemize}

\section{Why a Multivariate Distribution?}

Fit parameters are often correlated.

For example:

\begin{itemize}
    \item increasing one parameter,
    \item while decreasing another,
\end{itemize}

might produce similarly good fits.

The covariance matrix captures these correlations.

\section{Generating Many Curves}

Next:

\begin{lstlisting}[language=Python]
ysample = np.asarray(

    [

        fitfunction(xfit, *pi)

        for pi in ps

    ]

)
\end{lstlisting}

This creates:

\begin{itemize}
    \item 10,000 fitted curves,
    \item one for each sampled parameter set.
\end{itemize}

\section{Understanding the Shape of \texttt{ysample}}

The array:

\begin{lstlisting}[language=Python]
ysample
\end{lstlisting}

contains many different possible y-values for each x-position.

Each row corresponds to a different sampled fit.

\section{Constructing the Error Band}

We now compute percentiles.

Upper boundary:

\begin{lstlisting}[language=Python]
upper = np.percentile(
    ysample,
    84.0,
    axis=0
)
\end{lstlisting}

Lower boundary:

\begin{lstlisting}[language=Python]
lower = np.percentile(
    ysample,
    16.0,
    axis=0
)
\end{lstlisting}

\section{Why 16\% and 84\%?}

For a Gaussian distribution:

\begin{itemize}
    \item about 68\% of values lie within one standard deviation,
    \item leaving roughly 16\% in each tail.
\end{itemize}

Therefore:

\begin{itemize}
    \item the 16th percentile corresponds approximately to \(-1\sigma\),
    \item the 84th percentile corresponds approximately to \(+1\sigma\).
\end{itemize}

\section{Plotting the Band}

We can shade the region between the boundaries:

\begin{lstlisting}[language=Python]
plt.fill_between(
    xfit,
    lower,
    upper,
    alpha=0.3
)
\end{lstlisting}

This creates a shaded uncertainty region around the fit.

\section{Interpretation}

The error band gives a visual estimate of:

\begin{itemize}
    \item how uncertain the fit is,
    \item where predictions are more reliable,
    \item where uncertainties become large.
\end{itemize}

\section{Important Note}

This procedure assumes:

\begin{itemize}
    \item approximately Gaussian parameter uncertainties,
    \item a reasonable covariance matrix,
    \item a good fit model.
\end{itemize}

In professional data analysis, uncertainty estimation can become much more sophisticated.

\section{Complete Workflow}

Typical steps:

\begin{enumerate}
    \item Fit the data.
    \item Obtain \texttt{popt} and \texttt{pcov}.
    \item Generate many random parameter sets.
    \item Compute many fitted curves.
    \item Determine percentile boundaries.
    \item Plot the shaded region.
\end{enumerate}

\section{Key Takeaways}

\begin{itemize}
    \item Fit parameters have uncertainties.
    \item The covariance matrix describes parameter uncertainty and correlation.
    \item Monte Carlo sampling can propagate uncertainties into the fitted curve.
    \item Error bands provide a useful visualization of fit uncertainty.
\end{itemize}

\end{document}
